% ═══════════════════════════════════════════════════════════════
% LEHRERHINWEISE: Atommodelle (Physik Klasse 10)
% Stunde 1-2 (90 Minuten)
% ═══════════════════════════════════════════════════════════════

\documentclass[11pt, a4paper]{article}

\usepackage[utf8]{inputenc}
\usepackage[T1]{fontenc}
\usepackage[ngerman]{babel}
\usepackage{helvet}
\renewcommand{\familydefault}{\sfdefault}

\usepackage[margin=2cm]{geometry}
\usepackage{enumitem}
\usepackage{tabularx}
\usepackage{booktabs}
\usepackage{array}
\usepackage{xcolor}
\usepackage{tcolorbox}
\usepackage{fancyhdr}

\setlist{nosep, leftmargin=1.5em}

% Farben
\definecolor{physik}{HTML}{2c5aa0}
\definecolor{grau}{HTML}{888888}
\definecolor{hellgrau}{HTML}{f0f0f0}
\definecolor{warnung}{HTML}{b35c00}
\definecolor{warnungbg}{HTML}{fff3e0}

% Kopfzeile
\pagestyle{fancy}
\fancyhf{}
\renewcommand{\headrulewidth}{0.4pt}
\lhead{\small\textbf{LEHRERHINWEISE} -- Physik Klasse 10}
\rhead{\small Atommodelle (Stunde 1-2)}
\cfoot{\small\thepage}

% Boxen
\newtcolorbox{warnbox}[1][]{
    colback=warnungbg, colframe=warnung, boxrule=0.8pt, arc=0pt,
    left=6pt, right=6pt, top=6pt, bottom=6pt,
    fonttitle=\bfseries\small, title=#1
}

\newcommand{\abschnitt}[1]{%
    \vspace{10pt}\noindent\textbf{\textcolor{physik}{\large #1}}\vspace{6pt}\hrule\vspace{8pt}
}

\begin{document}

\begin{center}
{\LARGE\bfseries Lehrerhinweise: Atommodelle}\\[6pt]
{\large Physik Klasse 10 -- Stunde 1-2 (90 Min)}\\[4pt]
{\small\textcolor{grau}{Kernphysik-Einheit, Teil 1}}
\end{center}
\vspace{8pt}\hrule\vspace{12pt}

% ═══════════════════════════════════════════════════════════════
\abschnitt{1. Stundenverlauf (90 Minuten)}

\begin{tabularx}{\textwidth}{|c|c|X|l|}
\hline
\textbf{Zeit} & \textbf{Phase} & \textbf{Inhalt} & \textbf{Material} \\
\hline
0-5 & Einstieg & Frage: „Woraus besteht alles?" Brainstorming & Tafel \\
\hline
5-15 & Input & Lernpfad: Dalton → Thomson (historisch) & LP, Beamer \\
\hline
15-25 & Erarbeitung & AB Kasten 1 ausfüllen, Zeitstrahl & AB-01 \\
\hline
25-40 & Input & Lernpfad: Rutherford-Versuch (Video/Sim) & LP, PhET \\
\hline
40-50 & Erarbeitung & AB Kasten 2, Simulation beobachten & AB-01, SIM \\
\hline
50-60 & Input & Lernpfad: Bohr-Modell, Kernaufbau & LP \\
\hline
60-70 & Erarbeitung & AB Kasten 3, Übungen 1-3 & AB-01 \\
\hline
70-80 & Vertiefung & Übungen 4-5 (AFB III), Partnerarbeit & AB-01 \\
\hline
80-90 & Sicherung & Besprechung, Zusammenfassung, Ausblick & Tafel \\
\hline
\end{tabularx}

\vspace{12pt}

% ═══════════════════════════════════════════════════════════════
\abschnitt{2. Differenzierung}

\begin{tabularx}{\textwidth}{|c|X|X|}
\hline
\textbf{Niveau} & \textbf{Fokus} & \textbf{Hilfestellung} \\
\hline
★ (BR) & Kasten 1-3 ausfüllen, Übung 1 & Wortkarten für Lückentext, Partnerarbeit \\
\hline
★★ (MR) & Alle Kästen + Übungen 1-3 & Formelkarte für Nuklid-Schreibweise \\
\hline
★★★ (GY) & Komplettes AB inkl. Übung 4-5 & Selbstständig, Transfer-Aufgaben \\
\hline
\end{tabularx}

\vspace{12pt}

% ═══════════════════════════════════════════════════════════════
\abschnitt{3. Häufige Fehler \& Reaktionen}

\begin{warnbox}[Typische Schülerfehler]
\begin{enumerate}
\item \textbf{„Thomson wurde durch Elektronen widerlegt"}\\
→ Korrektur: Thomson entwickelte sein Modell WEGEN der Elektronenentdeckung. Rutherford widerlegte es.

\item \textbf{Verwechslung A und Z}\\
→ Eselsbrücke: \textbf{A} = \textbf{A}lle Nukleonen, \textbf{Z} = \textbf{Z}ahl der Protonen

\item \textbf{„Die meisten α-Teilchen wurden abgelenkt"}\\
→ Nachhaken: Was bedeutet das für das Atom? (Fast leer!) Nur 1:8000 stark abgelenkt.

\item \textbf{Größenverhältnis unterschätzt}\\
→ Vergleich: Wenn Kern = Stecknadelkopf (1mm), dann Atom = Fußballstadion (100m)
\end{enumerate}
\end{warnbox}

\vspace{12pt}

% ═══════════════════════════════════════════════════════════════
\abschnitt{4. Benötigtes Material}

\begin{itemize}
\item Lernpfad (Beamer/Tablets): \texttt{lernpfad-atommodelle.html}
\item Arbeitsblatt: \texttt{AB-01-atommodelle.pdf} (Klassensatz)
\item Musterlösung: \texttt{ML-01-atommodelle.pdf}
\item PhET-Simulation: \texttt{phet.colorado.edu/de/simulations/rutherford-scattering}
\item Optional: Periodensystem (Wandtafel oder Handout)
\end{itemize}

\vspace{12pt}

% ═══════════════════════════════════════════════════════════════
\abschnitt{5. Lösungen (Kurzfassung)}

\textbf{Kasten 1:} Dalton (1803, unteilbar), Thomson (1897, Rosinenkuchen), Rutherford (1911, Kern-Hülle), Bohr (1913, Bahnen)

\textbf{Kasten 2:} Alpha-Teilchen auf Goldfolie. Meiste durch → Atom leer. Wenige abgelenkt → Kern klein, positiv. Kern: $10^{-15}$m, Atom: $10^{-10}$m

\textbf{Kasten 3:} Proton (+1, 1u, Kern), Neutron (0, 1u, Kern), Elektron (-1, ≈0, Hülle). A=Massenzahl, Z=Ordnungszahl, N=A-Z

\textbf{Übung 3:} C-12: 6p, 6n, 6e / U-235: 92p, 143n, 92e / Fe-56: 26p, 30n, 26e

\vspace{8pt}
\textit{→ Vollständige Lösungen siehe ML-01-atommodelle.pdf}

\end{document}
